% 章节 01：条件期望最优估计与正交投影几何刻画
\section{最优估计}\label{sec:opt}

为具体说明所要研究的问题类型，考虑如下情形。设给定信号 $x_1(t)$ 与噪声 $x_2(t)$，而仅有二者的和 $y(t) = x_1(t) + x_2(t)$ 可供观测。假定我们已经观测并精确知晓 $y(t_0), \dots, y(t)$ 的取值。由此关于信号在时刻 $t = t_1$（其中 $t_1$ 可小于，等于或大于 $t$）处的不可直接观测取值，我们能够推断出什么？若 $t_1 < t$，这是一个\emph{数据平滑（内插）}问题；若 $t_1 = t$，称为\emph{滤波}；若 $t_1 > t$，则是一个\emph{预测}问题。由于本文的处理方式具有足够的一般性，可同时涵盖上述及类似问题，下文统一使用总括术语\emph{估计}。

% <!-- p.3 -->

正如 Wiener \cite{r01} 所指出的，估计问题的自然框架属于概率论与数理统计范畴。因此，信号，噪声及其和均为随机变量，随时间演化即构成随机过程。由随机过程的概率描述，我们可以确定信号与噪声的特定样本出现的概率。对于随机变量 $y(t)$ 的任意一组已测得数值 $\eta(t_0), \dots, \eta(t)$，原则上亦可确定随机变量 $x_1(t_1)$ 取各种可能实数值 $\xi_1$（已校订，见\hyperref[sec:k01]{校勘记 K-01}）的条件概率。这就是条件概率分布函数
\begin{equation}\label{eq:1}
\mathbb{P}\bigl(x_1(t_1) \leq \xi_1 \mid y(t_0) = \eta(t_0), \dots, y(t) = \eta(t)\bigr) = F(\xi_1)
\end{equation}

显然，$F(\xi_1)$ 集中体现了测量随机变量 $y(t_0), \dots, y(t)$ 所传递的关于随机变量 $x_1(t_1)$ 的全部信息。随机变量 $x_1(t_1)$ 的任何统计估计量都是该条件分布的某个泛函，从而是观测随机变量 $y(t_0), \dots, y(t)$ 的某个（非随机）函数。我们将这一统计估计量记作 $X_1(t_1 \mid t)$；当观测随机变量集或待估时刻由上下文不言自明时，亦简记作 $X_1(t_1)$ 或 $X_1$。

现设 $X_1$ 已给定为随机变量 $y(t_0), \dots, y(t)$ 的某个固定函数。则 $X_1$ 本身也是一个随机变量，且每当 $y(t_0), \dots, y(t)$ 的实际观测值已知时，$X_1$ 的实际取值便随之已知。一般而言，$X_1(t_1)$ 的实际取值不同于 $x_1(t_1)$ 的（未知）实际取值。为建立确定 $X_1$ 的合理准则，自然应当对不准确的估计赋予\emph{惩罚}或\emph{损失}。显然，该损失应当是关于\emph{估计误差} $\varepsilon = x_1(t_1) - X_1(t_1)$ 的 (i) 非负，(ii) 非减偶函数。据此，我们定义满足如下条件的\emph{损失函数}：
\begin{equation}\label{eq:2}
\begin{aligned}
&L(0) = 0 \\
&L(\varepsilon_2) \geq L(\varepsilon_1) \geq 0,\; \varepsilon_2 \geq \varepsilon_1 \geq 0 \\
&L(\varepsilon) = L(-\varepsilon)
\end{aligned}
\end{equation}

常见损失函数的例子有：$L(\varepsilon) = a\varepsilon^2$，$a\varepsilon^4$，$a|\varepsilon|$，$a\bigl[1 - \exp(-\varepsilon^2)\bigr]$ 等，其中 $a$ 为正常数。

选取随机变量 $X_1$ 的一种自然（但绝非唯一）方式，是要求该选取使平均损失（即风险）达到最小：
\begin{equation}\label{eq:3}
\mathbb{E}\bigl[L\bigl(x_1(t_1) - X_1(t_1)\bigr)\bigr] = \mathbb{E}\Bigl[\mathbb{E}\bigl[L\bigl(x_1(t_1) - X_1(t_1)\bigr) \mid y(t_0), \dots, y(t)\bigr]\Bigr]
\end{equation}

由于式~\eqref{eq:3} 右端的外层期望不依赖于 $X_1$ 的选取，而仅依赖于 $y(t_0), \dots, y(t)$，显然最小化式~\eqref{eq:3} 等价于对每组给定观测最小化条件期望损失
\begin{equation}\label{eq:4}
\mathbb{E}\bigl[L\bigl(x_1(t_1) - X_1(t_1)\bigr) \mid y(t_0), \dots, y(t)\bigr]
\end{equation}

仅需附加相当温和的假设，即可用简洁方式刻画最优估计。

% 定理 1：对称单峰条件分布下的最优估计等于条件数学期望
\begin{theorem}\label{thm:1}

设损失函数 $L$ 属于式~\eqref{eq:2} 所述类型，且由式~\eqref{eq:1} 定义的条件分布函数 $F(\xi)$ 满足：

(\emph{A}) 关于条件均值 $\bar{\xi}$ 对称（已校订对称中心，见\hyperref[sec:k02]{校勘记 K-02}）：
\begin{equation*}
F(\xi) = 1 - F(2\bar{\xi} - \xi)
\end{equation*}
（对含原子测度的一般分布，右端取左极限 $F\bigl((2\bar{\xi} - \xi)^-\bigr)$。）

(\emph{B}) 在区域 $\xi \leq \bar{\xi}$ 上为凸函数，即对任意 $\xi_1, \xi_2 \leq \bar{\xi}$ 及 $0 \leq \lambda \leq 1$ 有：
\begin{equation*}
F\bigl(\lambda\xi_1 + (1 - \lambda)\xi_2\bigr) \leq \lambda F(\xi_1) + (1 - \lambda)F(\xi_2)
\end{equation*}
则使平均损失 \eqref{eq:3} 达到最小的随机变量 $x_1^{\text{✫}}(t_1 \mid t)$ 恰为条件数学期望
\begin{equation}\label{eq:5}
x_1^{\text{✫}}(t_1 \mid t) = \mathbb{E}\bigl[x_1(t_1) \mid y(t_0), \dots, y(t)\bigr]
\end{equation}
\end{theorem}

\begin{proof}

正如 Sherman \cite{r25} 近期所指出的，本定理由概率论中一条熟知引理直接推出。

\end{proof}

\begin{corollary}\label{cor:1}

若随机过程 $\{x_1(t)\}$，$\{x_2(t)\}$ 与 $\{y(t)\}$ 为正态随机过程（原称 Gaussian random process），则定理~\ref{thm:1} 成立。

\end{corollary}

\begin{proof}

由\hyperref[sec:appendix]{附录}定理~\ref{thm:5} (\emph{A})，正态随机过程上的条件分布仍为正态分布，故定理~\ref{thm:1} 的各项要求恒得到满足。

\end{proof}

在控制系统文献中，这一结果有时以一种在某方面限制更严而在另一方面更具一般性的形式出现：

% 定理 1-a：二次损失下无分布对称性假设的最优估计
\begin{theoremn}{1-a}\label{thm:1a}

若 $L(\varepsilon) = \varepsilon^2$，则无需假设 (\emph{A}) 与 (\emph{B})，定理~\ref{thm:1} 的结论依然成立。

\end{theoremn}

\begin{proof}

将条件期望 \eqref{eq:4} 展开为
\begin{equation*}
\mathbb{E}\bigl[x_1^2(t_1) \mid y(t_0), \dots, y(t)\bigr] - 2X_1(t_1)\mathbb{E}\bigl[x_1(t_1) \mid y(t_0), \dots, y(t)\bigr] + X_1^2(t_1)
\end{equation*}
并对 $X_1(t_1)$ 求导置零即得。这一论证并非完全严格；简洁而严格的证明可参见 Doob \cite{r15} 第 77–78 页。

\end{proof}

\begin{remarkn}{(a)}\label{rem:a}

据作者所知，使得条件分布函数满足定理~\ref{thm:1} 要求的最一般随机过程类 $\{x_1(t)\}$，$\{x_2(t)\}$ 目前尚不清楚。

\end{remarkn}

\begin{remarkn}{(b)}\label{rem:b}

除 Sherman 的短文外，定理~\ref{thm:1} 似乎从未在控制系统文献中被明确陈述过。事实上，文献中常见到认为形如式~\eqref{eq:2} 的一般损失函数无法在数学上便捷处理的说法。

\end{remarkn}

\begin{remarkn}{(c)}\label{rem:c}

在下文中，我们将主要处理向量值随机变量。此时估计问题可表述为：给定向量值随机过程 $\{\bm{x}(t)\}$ 以及被观测随机向量 $\bm{y}(t_0), \dots, \bm{y}(t)$，其中 $\bm{y}(t) = \bm{M}\bm{x}(t)$（$\bm{M}$ 为列秩小于 $n$ 的矩阵，仅在嵌入为 $n \times n$ 方阵时称奇异矩阵，见\hyperref[sec:k03]{校勘记 K-03}；换言之，$\bm{x}(t)$ 的并非所有坐标分量均可直接观测），试求估计量 $\bm{X}(t_1)$ 以最小化期望损失 $\mathbb{E}\bigl[L\bigl(\|\bm{x}(t_1) - \bm{X}(t_1)\|\bigr)\bigr]$，其中 $\|\cdot\|$ 表示向量范数。

定理~\ref{thm:1} 在向量情形下依然成立，只要我们要求向量 $\bm{x}(t_1)$ 的 $n$ 个坐标分量在给定观测下的联合条件分布（其累积分布函数记为
\begin{equation*}
\mathbb{P}\bigl(x_1(t_1) \leq \xi_1, \dots, x_n(t_1) \leq \xi_n \mid \bm{y}(t_0), \dots, \bm{y}(t)\bigr) = F(\xi_1, \dots, \xi_n)
\end{equation*}
）关于条件均值向量 $\bar{\bm{\xi}} = (\bar{\xi}_1, \dots, \bar{\xi}_n)^{\mathrm{T}}$ 中心对称，且当存在条件概率密度 $p(\bm{\xi})$ 时其所有上水平集 $\{\bm{\xi} \in \mathbb{R}^n : p(\bm{\xi}) \geq c\}$ 均为凸集（即拟凹单峰条件密度；原文误作联合累积分布函数在负象限为凸函数，已校订，见\hyperref[sec:k04]{校勘记 K-04}）。

\end{remarkn}

% 第 4 节：随机变量子空间上的正交投影分解
\section{正交投影}\label{sec:proj}

将最优估计显式计算为被观测变量的函数，在一般情形下是不可能的。但有一个重要例外：随机过程 $\{x_1(t)\}$，$\{x_2(t)\}$ 为正态随机过程。

另一方面，若我们在二次损失限制 $L(\varepsilon) = \varepsilon^2$ 下寻求最优估计，并附加要求估计量是被观测随机变量的线性函数，则所得估计量与正态情形下无需线性限制或二次损失假设的最优估计完全相同。这表明，唯有当 (i) 随机过程为非正态过程，且即便如此（鉴于\hyperref[sec:appendix]{附录}定理~\ref{thm:5} (\emph{C})）亦唯有当 (ii) 至少计及三阶概率分布函数时，非线性估计才有可能改进线性估计所得的结果。

在上述特殊情形下，借助几何图像最容易理解估计问题的显式解。这正是本节的主题。

考虑实值随机变量 $y(t_0), \dots, y(t)$。这些随机变量以实数为系数的一切线性组合之集合
\begin{equation}\label{eq:6}
\sum_{i=t_0}^{t} a_i y(i)
\end{equation}
构成一个\emph{向量空间}——即\emph{线性子空间}（原文称线性流形），我们将其记作 $\mathcal{Y}(t)$。在抽象意义下，我们将形如式~\eqref{eq:6} 的任意表达式视为 $\mathcal{Y}(t)$ 中的一个“点”或“向量”；当然，此处“向量”一词的抽象用法不应与“向量值”随机变量等概念相混淆。由于我们不希望固定 $t$ 的取值（即可能观测的总次数），应将 $\mathcal{Y}(t)$ 视为一切可能观测所张成环境空间 $\mathcal{H}$ 的一个有限维线性子空间 $\mathcal{Y}(t) \leqslant \mathcal{H}$。

% <!-- p.4 -->

给定 $\mathcal{Y}(t)$ 中的任意两个向量 $u, v$（即形如式~\eqref{eq:6} 的随机变量），若 $\mathbb{E}[uv] = 0$，则称 $u$ 与 $v$ \emph{正交}。运用逐次正交归一化法（原称 Schmidt orthogonalization procedure，参见 Doob \cite{r15} 第 151 页或 Loève \cite{r16} 第 459 页），很容易在 $\mathcal{Y}(t)$ 中选取一组\emph{标准正交基}。这是指 $\mathcal{Y}(t)$ 中的一组向量 $e_{t_0}, \dots, e_t$（当观测线性相关时基元个数可相应减少），使得 $\mathcal{Y}(t)$ 中的任意向量均可唯一表示为 $e_{t_0}, \dots, e_t$ 的线性组合，且对 $i, j = t_0, \dots, t$ 满足
\begin{equation}\label{eq:7}
\begin{aligned}
&\mathbb{E}[e_i e_j] = \delta_{ij} = 1,\; i = j \\
&\mathbb{E}[e_i e_j] = \delta_{ij} = 0,\; i \neq j
\end{aligned}
\end{equation}
于是，$\mathcal{Y}(t)$ 中的任意向量 $\bar{x}$ 均可由下式给出：
\begin{equation*}
\bar{x} = \sum_{i=t_0}^{t} a_i e_i
\end{equation*}
从而各系数 $a_i$ 可借助式~\eqref{eq:7} 立即确定：
\begin{equation}\label{eq:8}
\mathbb{E}[\bar{x} e_j] = \mathbb{E}\left[\left(\sum_{i=t_0}^{t} a_i e_i\right) e_j\right] = \sum_{i=t_0}^{t} a_i \mathbb{E}[e_i e_j] = \sum_{i=t_0}^{t} a_i \delta_{ij} = a_j
\end{equation}

由此进一步推出，任意随机变量 $x$（未必属于 $\mathcal{Y}(t)$）均可唯一分解为两部分：属于 $\mathcal{Y}(t)$ 的部分 $\bar{x}$，以及正交于 $\mathcal{Y}(t)$（即正交于 $\mathcal{Y}(t)$ 中每一个向量）的部分 $\tilde{x}$。事实上，我们可写为
\begin{equation}\label{eq:9}
x = \bar{x} + \tilde{x} = \sum_{i=t_0}^{t} \mathbb{E}[x e_i]\, e_i + \tilde{x}
\end{equation}
因此 $\bar{x}$ 由方程~\eqref{eq:9} 唯一确定，且显然是 $\mathcal{Y}(t)$ 中的向量。从而 $\tilde{x}$ 亦被唯一确定；下面只需验证它正交于 $\mathcal{Y}(t)$：
\begin{equation*}
\mathbb{E}[\tilde{x} e_i] = \mathbb{E}[(x - \bar{x}) e_i] = \mathbb{E}[x e_i] - \mathbb{E}[\bar{x} e_i]
\end{equation*}

注意到 $\bar{x}$ 关于基底 $e_{t_0}, \dots, e_t$ 的坐标既可写成 $\mathbb{E}[\bar{x} e_i]$ 的形式（如式~\eqref{eq:8}），亦可写成 $\mathbb{E}[x e_i]$ 的形式（如式~\eqref{eq:9}）。由于基底展开坐标唯一，故对 $i = t_0, \dots, t$ 有 $\mathbb{E}[x e_i] = \mathbb{E}[\bar{x} e_i]$；由此得 $\mathbb{E}[\tilde{x} e_i] = 0$，即 $\tilde{x}$ 正交于每一个基向量 $e_i$，从而正交于 $\mathcal{Y}(t)$。我们称 $\bar{x}$ 为 $x$ 在 $\mathcal{Y}(t)$ 上的\emph{正交投影}。

正交投影还可用另一种等价方式刻画：$\bar{x}$ 是 $\mathcal{Y}(t)$ 中（即随机变量 $y(t_0), \dots, y(t)$ 的一切\emph{线性}函数中）使二次损失函数达到最小的那个向量。事实上，若 $\bar{w}$ 为 $\mathcal{Y}(t)$ 中的任意其他向量，则有
\begin{equation*}
\mathbb{E}[(x - \bar{w})^2] = \mathbb{E}[(\tilde{x} + \bar{x} - \bar{w})^2] = \mathbb{E}\bigl[\bigl((x - \bar{x}) + (\bar{x} - \bar{w})\bigr)^2\bigr]
\end{equation*}
由于 $\tilde{x}$ 正交于 $\mathcal{Y}(t)$ 中的每一个向量，特别地正交于 $\bar{x} - \bar{w}$，交叉项为零，故有
\begin{equation}\label{eq:10}
\mathbb{E}[(x - \bar{w})^2] = \mathbb{E}[(x - \bar{x})^2] + \mathbb{E}[(\bar{x} - \bar{w})^2] \geq \mathbb{E}[(x - \bar{x})^2]
\end{equation}
这表明，若 $\bar{w}$ 同样使二次损失达到最小，则必有 $\mathbb{E}[(\bar{x} - \bar{w})^2] = 0$，这意味着随机变量 $\bar{x}$ 与 $\bar{w}$ 相等（至多除去一个概率为零的事件集）。

上述结果可归纳如下：

% 定理 2：正态过程或线性二次估计下的正交投影等价性
\begin{theorem}\label{thm:2}

设 $\{x(t)\}$，$\{y(t)\}$ 为零均值随机过程（即对一切 $t$ 有 $\mathbb{E}[x(t)] = \mathbb{E}[y(t)] = 0$），且观测到 $y(t_0), \dots, y(t)$。若满足以下两条件之一：

(\emph{A}) 随机过程 $\{x(t)\}$，$\{y(t)\}$ 为联合正态随机过程（且损失函数 $L$ 属于式~\eqref{eq:2} 所述类型）；或

(\emph{B}) 最优估计限定为被观测随机变量的线性函数，且 $L(\varepsilon) = \varepsilon^2$；

则有
\begin{equation}\label{eq:11}
\begin{aligned}
&x^{\text{✫}}(t_1 \mid t) = \text{给定 } y(t_0), \dots, y(t) \text{ 时 } x(t_1) \text{ 的最优估计} \\
= &x(t_1) \text{ 在 } \mathcal{Y}(t) \text{ 上的正交投影 } \bar{x}(t_1 \mid t)
\end{aligned}
\end{equation}
\end{theorem}

这些结果在概率论中已为人熟知，但在控制系统文献中尚不易查得；参见 Doob \cite{r15} 第 75–78 页或 Pugachev \cite{r26}。为便于算子化表述，我们常将向子空间 $\mathcal{Y}(t)$ 的正交投影算子与余投影算子分别记作
\begin{equation*}
\begin{aligned}
&\bar{x}(t_1 \mid t) \equiv x^{\text{✫}}(t_1 \mid t) = \mathscr{P}_{\mathcal{Y}(t)}[x(t_1)] \\
&\tilde{x}(t_1 \mid t) = x(t_1) - \bar{x}(t_1 \mid t) = \mathscr{P}_{\mathcal{Y}(t)}^{\perp}[x(t_1)]
\end{aligned}
\end{equation*}

按\hyperref[thm:2]{定理}第 (\emph{A}) 部分（原文误引作 (\emph{b})，并颠倒了证明段落标号，已一并校订，见\hyperref[sec:k05]{校勘记 K-05}）：当所论随机过程为正态随机过程时，正交投影在数学上与条件期望完全恒等。

\begin{proof}

(\emph{A}) 由于 $x(t)$ 与 $y(t)$ 均为零均值随机变量，由公式 \eqref{eq:9} 显然可知，$x(t_1)$ 关于线性子空间 $\mathcal{Y}(t)$ 的正交余部 $\tilde{x}(t_1 \mid t)$ 亦为零均值随机变量。零均值正交随机变量必不相关；若它们同时又是联合正态随机变量，则由\hyperref[sec:appendix]{附录}定理~\ref{thm:5} (\emph{B}) 知它们相互独立。因此
\begin{equation*}
\begin{aligned}
&0 = \mathbb{E}[\tilde{x}(t_1 \mid t)] = \mathbb{E}\bigl[\tilde{x}(t_1 \mid t) \mid y(t_0), \dots, y(t)\bigr] \\
= &\mathbb{E}\bigl[x(t_1) - \bar{x}(t_1 \mid t) \mid y(t_0), \dots, y(t)\bigr] \\
= &\mathbb{E}\bigl[x(t_1) \mid y(t_0), \dots, y(t)\bigr] - \bar{x}(t_1 \mid t) = 0
\end{aligned}
\end{equation*}
结合定理~\ref{thm:1} \hyperref[cor:1]{推论}，即知 $\bar{x}(t_1 \mid t)$ 等于条件期望从而为最优估计。

(\emph{B}) 这直接由关于式~\eqref{eq:10} 的讨论推出。

\end{proof}

\begin{remarkn}{(d)}\label{rem:d}

当 $t \to \infty$ 时，要对本节内容给出严格的数学表述，需要用到完备内积空间（原称 Hilbert space）理论中的若干基本概念；参见 Doob \cite{r15} 与 Loève \cite{r16}。

\end{remarkn}

\begin{remarkn}{(e)}\label{rem:e}

对定理~\ref{thm:2} 的物理解释在很大程度上取决于研究者的取向。若我们对正态性假设不抱顾虑，则第 (\emph{A}) 部分表明，对于一切合理的损失函数，正交投影均是最优估计。反之，若我们对正态性存有顾虑，则即便退而只考虑线性估计，我们也知道对于许多合理的损失函数，正交投影\emph{并非}最优估计。由于在实际问题中很难确切判定一个物理起源的随机过程在何种近似程度上具有正态性，因此很难断言定理~\ref{thm:2} 究竟具有极其广泛的意义还是相当有限的意义。

\end{remarkn}

\begin{remarkn}{(f)}\label{rem:f}

定理~\ref{thm:2} 可立即推广至向量值随机变量的情形。事实上，我们将由随机向量 $\bm{y}(t_0), \dots, \bm{y}(t)$ 张成的线性子空间 $\mathcal{Y}(t)$ 定义为各随机向量 $\bm{y}(t_0), \dots, \bm{y}(t)$ 的全部 $m$ 个坐标分量的一切实系数线性组合之集合：
\begin{equation*}
\sum_{i=t_0}^{t} \sum_{j=1}^{m} a_{ij} y_j(i)
\end{equation*}

其余推导过程与标量情形如出一辙（对随机向量的正交投影算子 $\mathscr{P}_{\mathcal{Y}(t)}[\bm{x}(t_1)]$ 按各坐标分量逐分量向 $\mathcal{Y}(t)$ 投影理解）。

\end{remarkn}

\begin{remarkn}{(g)}\label{rem:g}

定理~\ref{thm:2} 实质上表明，在条件 (\emph{A}) 或 (\emph{B}) 下，最优估计是此前全部观测值的线性组合。换言之，最优估计可视为一个线性滤波器的输出，其输入为可观测随机变量实际发生的样本值；定理~\ref{thm:2} 给出了一种计算该最优滤波器单位冲激响应的途径。正如前文所指出的，仅知晓这一单位冲激响应并不构成问题的完整解答；基于这一原因，本文不再给出计算单位冲激响应的显式公式。

\end{remarkn}
